\documentclass[aspectratio=169,11pt]{beamer}
% Shared Beamer style for practice contest solution walkthroughs.
\usepackage[utf8]{inputenc}
\usepackage[T1]{fontenc}
\usepackage{lmodern}
\usepackage{amsmath,amssymb}
\usepackage{xcolor}
\usepackage{hyperref}
\usepackage{listings}

\definecolor{CPNavy}{HTML}{172033}
\definecolor{CPBlue}{HTML}{2563EB}
\definecolor{CPGreen}{HTML}{16A34A}
\definecolor{CPOrange}{HTML}{F97316}
\definecolor{CPGray}{HTML}{64748B}
\definecolor{CPLight}{HTML}{F8FAFC}
\definecolor{CPCodeBg}{HTML}{F1F5F9}
\definecolor{CPCodeRule}{HTML}{CBD5E1}

\usetheme{default}
\usefonttheme{professionalfonts}
\setbeamertemplate{navigation symbols}{}
\setbeamersize{text margin left=0.65cm,text margin right=0.65cm}
\setbeamertemplate{itemize item}{\color{CPBlue}$\blacktriangleright$}
\setbeamertemplate{itemize subitem}{\color{CPGray}$\triangleright$}

\setbeamercolor{normal text}{fg=CPNavy,bg=white}
\setbeamercolor{frametitle}{fg=white,bg=CPNavy}
\setbeamercolor{title}{fg=white,bg=CPNavy}
\setbeamercolor{subtitle}{fg=CPLight,bg=CPNavy}
\hypersetup{colorlinks=true,urlcolor=CPBlue}

\setbeamertemplate{frametitle}{%
  \nointerlineskip%
  \begin{beamercolorbox}[wd=\paperwidth,ht=0.88cm,dp=0.22cm,leftskip=0.55cm,rightskip=0.35cm]{frametitle}
    \usebeamerfont{frametitle}\insertframetitle\hfill\small\insertframenumber/\inserttotalframenumber
  \end{beamercolorbox}%
}

\setbeamertemplate{footline}{%
  \leavevmode%
  \hbox{%
    \begin{beamercolorbox}[wd=.58\paperwidth,ht=2.4ex,dp=1ex,leftskip=.5cm]{author in head/foot}%
      \color{CPGray}\insertshorttitle
    \end{beamercolorbox}%
    \begin{beamercolorbox}[wd=.42\paperwidth,ht=2.4ex,dp=1ex,rightskip=.5cm plus1fil]{date in head/foot}%
      \hfill\color{CPGray}\insertshortauthor
    \end{beamercolorbox}%
  }%
}

\setbeamertemplate{title page}{%
  \vspace*{1.25cm}
  \begin{beamercolorbox}[wd=\paperwidth,sep=0.65cm,leftskip=0.15cm,rightskip=0.15cm]{title}
    {\color{white}\Huge\bfseries\inserttitle\par}
    \vspace{0.35cm}
    {\color{CPLight}\Large\insertsubtitle\par}
    \vspace{0.6cm}
    {\color{CPLight}\large\insertauthor\par}
  \end{beamercolorbox}
  \vfill
}

\newcommand{\sectiontag}[1]{\textbf{\color{CPBlue}#1}}
\newenvironment{tightitemize}{\begin{itemize}\small\setlength{\itemsep}{0.18cm}}{\end{itemize}}
\lstdefinestyle{cpcode}{
  language=Python,
  basicstyle=\ttfamily\tiny,
  keywordstyle=\color{CPBlue}\bfseries,
  commentstyle=\color{CPGray}\itshape,
  stringstyle=\color{CPGreen},
  backgroundcolor=\color{CPCodeBg},
  frame=single,
  framerule=0.25pt,
  rulecolor=\color{CPCodeRule},
  showstringspaces=false,
  columns=fullflexible,
  keepspaces=true,
  breaklines=true,
  breakatwhitespace=false,
  tabsize=4,
  xleftmargin=0.08cm,
  xrightmargin=0.08cm,
  aboveskip=0.08cm,
  belowskip=0.08cm
}


\title[knapsack]{Knapsack}
\subtitle{Day 3 solution walkthrough}
\author[Solution Deck]{Competitive Programming Bootcamp}
\date{}

\begin{document}

\begin{frame}[plain]
  \titlepage
\end{frame}

% BEGIN GENERATED STATEMENT INTERPRETATION
\begin{frame}{Interpret the Word Problem}
\small
\sectiontag{Statement excerpts:}
\begin{tightitemize}
  \item \textit{``greatest value that fit into the knapsack''}
  \item \textit{``indices of the chosen items''}
\end{tightitemize}

\vspace{0.18cm}
\sectiontag{Programming interpretation:}
\begin{tightitemize}
  \item This is 0/1 knapsack: each item can be taken at most once.
  \item The score is value, while the constraint is total weight at most capacity.
  \item Because output requires item indices, the DP must support reconstruction.
\end{tightitemize}
\end{frame}
% END GENERATED STATEMENT INTERPRETATION







\begin{frame}{Problem Model}
\small
\sectiontag{Textbook connection:} CP4 3.5 c. Knapsack / Subset-Sum

\vspace{0.25cm}
\sectiontag{Core technique:} 0/1 knapsack DP with reconstruction

\vspace{0.25cm}
\begin{tightitemize}
  \item Each item has a value and weight.
  \item The capacity limits total selected weight.
  \item The output asks for selected item indices, not just the best value.
\end{tightitemize}
\end{frame}

\begin{frame}{How to Recognize the Approach}
\begin{tightitemize}
  \item Each item can be taken at most once, so this is 0/1 knapsack.
  \item A one-dimensional dp array gives best value per capacity.
  \item A chosen table is needed to reconstruct the selected indices.
\end{tightitemize}
\end{frame}

\begin{frame}{Algorithm}
\begin{tightitemize}
  \item Initialize dp[w] = 0 for every capacity w.
  \item For each item i, scan capacities backward from capacity to weight.
  \item If dp[w - weight] + value improves dp[w], update dp[w].
  \item Mark chosen[i][w] = 1 for improved states.
  \item Reconstruct by walking items backward and subtracting chosen item weights.
\end{tightitemize}
\end{frame}

\begin{frame}{Walkthrough of the Key State}
\begin{tightitemize}
  \item Backward capacity order prevents using item i more than once.
  \item chosen[i][w] means item i was responsible for the best known value at capacity w.
  \item The reconstruction follows those choices from the final capacity.
\end{tightitemize}
\end{frame}

\begin{frame}{Why the Algorithm Is Correct}
\begin{tightitemize}
  \item After processing i items, dp[w] is the best value using only those items within capacity w.
  \item The transition considers both excluding and including the current item.
  \item The chosen table records exactly the include decisions that produced the final DP value.
\end{tightitemize}
\end{frame}

\begin{frame}{Complexity and Implementation Notes}
\small
\sectiontag{Complexity:} O(NC) time and O(NC) memory because reconstruction stores choices for every item/capacity.

\vspace{0.3cm}
\begin{tightitemize}
  \item Multiple test cases are read until EOF.
  \item If only the value were required, memory could be reduced further.
\end{tightitemize}
\end{frame}



% BEGIN GENERATED CODE WALKTHROUGH
\begin{frame}[fragile]{Code: Initialize DP Tables}
\scriptsize
\sectiontag{Source lines:} \texttt{knapsack.py:33--42}

\vspace{0.08cm}
\begin{columns}[T,totalwidth=\textwidth]
\begin{column}{0.58\textwidth}
\begin{lstlisting}[style=cpcode]
import sys

def solve(capacity, n, items):
    dp = [0] * (capacity + 1)

    chosen = [bytearray(capacity + 1) for _ in range(n)]
\end{lstlisting}
\end{column}
\begin{column}{0.38\textwidth}
\sectiontag{What this chunk does}
\begin{tightitemize}
  \item dp[w] stores the best value available with capacity w.
  \item chosen[i][w] remembers that item i improved capacity w.
  \item The bytearray table keeps reconstruction memory compact.
\end{tightitemize}
\end{column}
\end{columns}
\end{frame}

\begin{frame}[fragile]{Code: Process Each Item Once}
\scriptsize
\sectiontag{Source lines:} \texttt{knapsack.py:44--53}

\vspace{0.08cm}
\begin{columns}[T,totalwidth=\textwidth]
\begin{column}{0.58\textwidth}
\begin{lstlisting}[style=cpcode]
for i in range(n):
    value, weight = items[i]

    for w in range(capacity, weight - 1, -1):
        new_value = dp[w - weight] + value

        if new_value > dp[w]:
            dp[w] = new_value
            chosen[i][w] = 1
\end{lstlisting}
\end{column}
\begin{column}{0.38\textwidth}
\sectiontag{What this chunk does}
\begin{tightitemize}
  \item For each item, capacities are scanned backward.
  \item Backward scanning prevents reusing the same item multiple times.
  \item When taking the item improves value, both dp and chosen are updated.
\end{tightitemize}
\end{column}
\end{columns}
\end{frame}

\begin{frame}[fragile]{Code: Reconstruct Chosen Items}
\scriptsize
\sectiontag{Source lines:} \texttt{knapsack.py:55--66}

\vspace{0.08cm}
\begin{columns}[T,totalwidth=\textwidth]
\begin{column}{0.58\textwidth}
\begin{lstlisting}[style=cpcode]
result = []
remaining = capacity

for i in range(n - 1, -1, -1):
    if chosen[i][remaining]:
        result.append(i)
        remaining -= items[i][1]

result.reverse()

return result
\end{lstlisting}
\end{column}
\begin{column}{0.38\textwidth}
\sectiontag{What this chunk does}
\begin{tightitemize}
  \item Reconstruction starts at full capacity after all items are processed.
  \item If chosen[i][remaining] is set, item i was used in the optimal path.
  \item The result is reversed because reconstruction walks from last item to first.
\end{tightitemize}
\end{column}
\end{columns}
\end{frame}

\begin{frame}[fragile]{Code: Read Multiple Cases}
\scriptsize
\sectiontag{Source lines:} \texttt{knapsack.py:69--89}

\vspace{0.08cm}
\begin{columns}[T,totalwidth=\textwidth]
\begin{column}{0.58\textwidth}
\begin{lstlisting}[style=cpcode]
def main():
    input = sys.stdin.buffer.readline
    output = []

    while True:
        line = input()

        if not line:
            break

        if not line.strip():
            continue

        capacity, n = map(int, line.split())

        items = []

        for _ in range(n):
            value, weight = map(int, input().split())
            items.append((value, weight))
\end{lstlisting}
\end{column}
\begin{column}{0.38\textwidth}
\sectiontag{What this chunk does}
\begin{tightitemize}
  \item The input has an unknown number of test cases, so the loop stops at EOF.
  \item Blank lines are skipped defensively.
  \item Each case reads the capacity and all value/weight pairs.
\end{tightitemize}
\end{column}
\end{columns}
\end{frame}

\begin{frame}[fragile]{Code: Output Reconstructed Items}
\scriptsize
\sectiontag{Source lines:} \texttt{knapsack.py:91--102}

\vspace{0.08cm}
\begin{columns}[T,totalwidth=\textwidth]
\begin{column}{0.58\textwidth}
\begin{lstlisting}[style=cpcode]
        result = solve(capacity, n, items)

        output.append(str(len(result)))

        output.append(" ".join(map(str, result)))

    sys.stdout.write("\n".join(output))

main()
\end{lstlisting}
\end{column}
\begin{column}{0.38\textwidth}
\sectiontag{What this chunk does}
\begin{tightitemize}
  \item solve returns the chosen item indices for the case.
  \item The final output follows Kattis: count first, then selected indices.
  \item All output lines are buffered and written once.
\end{tightitemize}
\end{column}
\end{columns}
\end{frame}
% END GENERATED CODE WALKTHROUGH

\begin{frame}{Source}
\small
\sectiontag{Solution file:} \texttt{practice\_contests/day\_3/knapsack.py}

\vspace{0.3cm}
\sectiontag{Problem:} \url{https://open.kattis.com/problems/knapsack}

\vspace{0.45cm}
Use this deck with the source code open beside it. The slides explain the reasoning; the code shows the exact input handling and output formatting.
\end{frame}

\end{document}
